\documentclass[10pt]{article}
\usepackage[margin=1in]{geometry}
\usepackage{amsmath,amssymb,amsthm}
\usepackage[colorlinks=true,urlcolor=blue,linkcolor=blue]{hyperref}
\newtheorem{theorem}{Theorem}
\newtheorem{lemma}{Lemma}
\newcommand{\rf}[2]{(#1)_{#2}}
\title{Positive coefficients of rising-factorial Tur\'an polynomials\\
\large A proof of the Karp--Sitnik conjecture, MathDB 323156}
\author{Principia Math In Progress}
\date{11 September 2026}
\begin{document}
\maketitle

\section*{Task 1. Prove coefficient positivity}

Write $(x)_k=x(x+1)\cdots(x+k-1)$, with $(x)_0=1$.
For a nonnegative sequence $(f_k)_{k=0}^n$ and $\alpha,\beta>0$, define
\[
 Q_n^{\alpha,\beta}(x)=
 \sum_{k=0}^n \binom nk f_kf_{n-k}
 \bigl[(x+\alpha)_k(x+\beta)_{n-k}
       -(x+\alpha+\beta)_k(x)_{n-k}\bigr].
\]

\begin{theorem}
If $n\geq3$ and $f_k^2>f_{k-1}f_{k+1}$ for $1\leq k\leq n-1$,
then $Q_n^{\alpha,\beta}$ has degree $n-2$ and a strictly positive
coefficient at every $x^j$, $0\leq j\leq n-2$, for all $\alpha,\beta>0$.
\end{theorem}

This is Conjecture 1 in Karp~\cite{K}. The proof first treats central
truncations of the rising-factorial Vandermonde sum, and then uses
summation by parts to handle the sequence weights.

\section*{A positive derivative quotient}

For integers $1\leq r\leq\lfloor n/2\rfloor$, set
\[
 C_{n,r}(A,B)=
 \sum_{k=r}^{n-r}\binom nk(A)_k(B)_{n-k}.
\]
This is a symmetric polynomial in $A,B$.

\begin{lemma}\label{lem:derivative}
There is a polynomial $R_{n,r}(A,B)$ with nonnegative coefficients,
of degree exactly $n-2$, such that
\begin{equation}
 (\partial_A-\partial_B)C_{n,r}(A,B)
 =(B-A)R_{n,r}(A,B).
 \label{eq:derivative}
\end{equation}
\end{lemma}

\begin{proof}
The elementary identity
\begin{equation}
 \frac{d}{dA}(A)_k
 =\sum_{s=1}^k\binom ks(s-1)!(A)_{k-s}
 \label{eq:risingderivative}
\end{equation}
follows by differentiating
$\sum_{k\geq0}(A)_kz^k/k!=(1-z)^{-A}$ and using
$-\log(1-z)=\sum_{s\geq1}z^s/s$.
All subsequent computations are identities of finite polynomials.

Apply~\eqref{eq:risingderivative} to $C_{n,r}$, and fix $s$.
Put $N=n-s$. In $\partial_A C_{n,r}$ the remaining first index ranges
over $[r-s,N-r]$; in $\partial_B C_{n,r}$ it ranges over
$[r,n-r]$. Both intervals are intersected with $[0,N]$.
They are reflections of each other under $k\mapsto N-k$.
Their common terms cancel. With
\[
 L=\max(0,r-s),\qquad U=\min(r-1,N-r),
\]
the surviving difference is therefore
\begin{equation}
 \begin{split}
 (\partial_A-\partial_B)C_{n,r}
 ={}&\sum_{s=1}^{n-r}\binom ns(s-1)!
 \sum_{k=L}^{U}\binom Nk\\
 &\quad\cdot\bigl[(A)_k(B)_{N-k}
                    -(A)_{N-k}(B)_k\bigr].
 \end{split}
 \label{eq:boundary}
\end{equation}
An empty inner sum is interpreted as zero.
For every surviving $k$, the bounds $k\leq r-1$ and $k\leq N-r$
give $2k\leq N-1$.
Hence $m=N-2k\geq1$, and the bracket factors as
\[
 (A)_k(B)_k\bigl[(B+k)_m-(A+k)_m\bigr].
\]
The polynomial $P(t)=(t+k)_m$ has nonnegative coefficients.
If $P(t)=\sum_i p_i t^i$, its divided difference is
\begin{equation}
 \frac{P(B)-P(A)}{B-A}
 =\sum_{i\geq1}p_i\sum_{h=0}^{i-1}A^hB^{i-1-h},
 \label{eq:quotient}
\end{equation}
which has nonnegative coefficients. This is a polynomial identity,
including $A=B$.
Dividing~\eqref{eq:boundary} by $B-A$ thus constructs the required
polynomial $R_{n,r}$ with nonnegative coefficients.

Each summand has degree at most $N-1=n-s-1\leq n-2$.
For $s=1$ the inner sum contains exactly $k=r-1$; then
$m=n-2r+1\geq1$, and its degree is
$2(r-1)+(m-1)=n-2$ with positive leading coefficients.
No cancellation is possible. Thus $\deg R_{n,r}=n-2$.
\end{proof}

\section*{Balancing the two parameters}

\begin{lemma}\label{lem:central}
For every $\alpha,\beta>0$, the polynomial
\[
 V_{n,r}(x)=C_{n,r}(x+\alpha,x+\beta)
            -C_{n,r}(x+\alpha+\beta,x)
\]
has strictly positive coefficients at all powers $x^j$,
$0\leq j\leq n-2$, and degree $n-2$.
\end{lemma}

\begin{proof}
Put $c=\alpha+\beta$ and $h=\min(\alpha,\beta)$.
Along $A=x+t$, $B=x+c-t$, equation~\eqref{eq:derivative} gives
\[
 \frac{d}{dt}C_{n,r}(x+t,x+c-t)
 =(c-2t)R_{n,r}(x+t,x+c-t).
\]
Using symmetry at the two endpoints, integrate to obtain
\begin{equation}
 V_{n,r}(x)=
 \int_0^h(c-2t)R_{n,r}(x+t,x+c-t)\,dt.
 \label{eq:integral}
\end{equation}
This identity may be integrated coefficient by coefficient since the
integrand is a finite polynomial in $x,t$.

For $0<t<h$, both $t$ and $c-t$ are positive, and $c-2t>0$.
All coefficients of $R_{n,r}$ are nonnegative. Furthermore it has a
monomial $A^uB^v$ with positive coefficient and $u+v=n-2$.
The polynomial $(x+t)^u(x+c-t)^v$ has strictly positive coefficients
at every degree from $0$ to $n-2$.
Thus every required coefficient of the integrand is strictly positive
throughout $(0,h)$. Equation~\eqref{eq:integral} proves the lemma.
This argument includes $\alpha=\beta$: only the integrand's endpoint
factor then vanishes.
\end{proof}

\section*{The sequence weights}

Let $M=\lfloor n/2\rfloor$ and $w_k=f_kf_{n-k}$.
The weights are symmetric and
\begin{equation}
 w_0<w_1<\cdots<w_M.
 \label{eq:weights}
\end{equation}
Indeed the strict log-concavity assumption implies $f_k>0$ for
$1\leq k\leq n-1$.
When both endpoints are positive, the successive ratios
$f_k/f_{k-1}$ strictly decrease; for $1\leq k\leq M$ this gives
\[
 \frac{w_k}{w_{k-1}}
 =\frac{f_k/f_{k-1}}{f_{n-k+1}/f_{n-k}}>1.
\]
If an endpoint is zero, $w_0=0<w_1$ and the same ratio argument applies
to all $k\geq2$, using only positive interior terms.

Define $C_{n,0}$ and $V_{n,0}$ by the same sums with $r=0$.
Vandermonde's identity gives $C_{n,0}(A,B)=(A+B)_n$, so
$V_{n,0}(x)=0$.
Writing
$w_k=w_0+\sum_{r=1}^{\min(k,n-k)}(w_r-w_{r-1})$
and interchanging finite sums now gives the exact identity
\begin{equation}
 Q_n^{\alpha,\beta}(x)
 =\sum_{r=1}^{M}(w_r-w_{r-1})V_{n,r}(x).
 \label{eq:decomposition}
\end{equation}
Every multiplier is positive by~\eqref{eq:weights}, and every
$V_{n,r}$ has precisely the strict coefficient positivity required
by Lemma~\ref{lem:central}. This proves the theorem.

The argument in fact only needs symmetric weights nondecreasing
toward the center, with at least one positive increment.
It is valid for $n=2$ as well, where the conclusion is a positive constant.

\section*{Verification and source status}

The proof is a finite algebraic argument and does not depend on a numerical
search. The accompanying verifier expands the original central sums
independently of the constructive positive quotient, checks their
derivative identity exactly, and checks the weighted decomposition.
The full published statement and current public claim sources were
checked on 11 September 2026; no matching full solution claim was located.
The source's known $\alpha=\beta=1$ theorem is distinguished from the
unrestricted shifts proved here in \texttt{literature-audit.md}.

\section*{Tasks 2--5}

\begin{enumerate}
\setcounter{enumi}{1}
\item Characterize all symmetric weight vectors for which the weighted
Turan polynomial has nonnegative coefficients for every $\alpha,\beta>0$,
and determine how far the sufficient monotonicity condition can be weakened.
\item Determine the sharp common translation range $x\mapsto x-\delta$
that preserves coefficient positivity for all admissible sequences,
as a function of $\alpha,\beta$ and $n$.
\item Prove or refute the companion real-negative-zero conjecture for
$Q_n^{1,1}(x-1)$ when the input sequence is a P\'olya frequency sequence
of infinite order.
\item Extend the positive derivative-quotient method to higher-order
Toeplitz determinants of rising-factorial series and determine sequence
conditions that ensure positive coefficients.
\end{enumerate}

\begin{thebibliography}{9}
\bibitem{K}
D.~Karp, \emph{Positivity of Toeplitz determinants formed by rising
factorial series and properties of related polynomials},
Zap. Nauchn. Sem. POMI \textbf{404} (2012), 184--198,
Conjecture 1 and formula (3), p.~185.
\href{https://www.mathnet.ru/eng/znsl5267}{Published article};
\href{https://arxiv.org/abs/1203.1482}{arXiv:1203.1482}.
English translation: Journal of Mathematical Sciences
\textbf{193} (2013), 106--114,
\href{https://doi.org/10.1007/s10958-013-1438-y}
{doi:10.1007/s10958-013-1438-y}.
\end{thebibliography}
\end{document}
